% ============================================================================
% CAPÍTULO 15 — Full stack\index{full stack} integrado com Next.js
% Status: texto completo (rodada editorial 2)
% ============================================================================
\chapter{Full Stack Integrado com Next.js}
\label{cap:fullstack}

\begin{caixaconceito}
Aqui tudo se conecta: React (capítulo~\ref{cap:react}), Next.js
(capítulo~\ref{cap:nextjs}), PostgreSQL (capítulo~\ref{cap:postgresql}),
Prisma\index{Prisma} (capítulo~\ref{cap:orms}) e autenticação (capítulo~\ref{cap:auth})
formando um \textbf{produto full stack único} — sem servidores separados para
manter. Server Actions\index{Server Actions}, Route Handlers e cache sob demanda são o coração deste
capítulo.
\end{caixaconceito}

Até aqui, você aprendeu as peças separadamente: uma API REST com Express ou
Hono, um banco relacional com Prisma, um frontend React. Agora o modelo muda de
figura. No App Router do Next.js, \emph{frontend e backend vivem no mesmo
projeto} — às vezes, no mesmo arquivo. Isso não é uma moda: é a evolução
natural de uma década de separação forçada entre cliente e servidor, e é o
modelo que o mercado de TypeScript mais contratou nos últimos anos.

Ao final deste capítulo, você será capaz de:
\begin{objetivos}
  \item Conectar o Next.js ao PostgreSQL via Prisma em Server Components;
  \item Implementar mutações com Server Actions e \texttt{useActionState};
  \item Criar APIs com Route Handlers quando necessário;
  \item Invalidar cache sob demanda com \texttt{revalidatePath}/\texttt{revalidateTag};
  \item Proteger páginas e ações com autenticação e RBAC;
  \item Entregar o projeto do capítulo: um marketplace de serviços.
\end{objetivos}

\section{A arquitetura: uma única aplicação, três camadas}

No modelo clássico, você mantinha pelo menos dois projetos: um \texttt{api/}
(backend) e um \texttt{web/} (frontend), conversando por HTTP em produção.
Isso funciona, mas paga um custo: dois deploys, dois ambientes, duas formas de
validar dados, e uma camada de rede — com latência, serialização e segurança —
entre código que poderia simplesmente conversar na memória.

O App Router propõe algo diferente. No mesmo projeto, quatro tipos de código
convivem, cada um com um papel claro:

\begin{itemize}[leftmargin=*, itemsep=0.4em]
  \item \textbf{Server Components} que buscam dados direto do banco (sem camada de API intermediária);
  \item \textbf{Server Actions} que escrevem no banco a partir de formulários (mutações com \texttt{<form action=\{\ldots\}>});
  \item \textbf{Route Handlers} (\texttt{route.ts}) para APIs públicas ou integrações;
  \item \textbf{Client Components} apenas para interatividade pontual.
\end{itemize}

A Figura~\ref{fig:skillhub-arquitetura} mostra como essas camadas se
organizam no SkillHub: o navegador fala apenas com o Next.js, que acessa o
banco (síncrono) e as filas (assíncrono) — sem uma API intermediária
separada.

\begin{figure}[htbp]
  \centering
  \diagramatikz{%
  \begin{tikzpicture}[
      box/.style={draw=primaria, fill=azulclaro, rounded corners=2pt,
        align=center, inner sep=4pt, font=\footnotesize},
      seta/.style={-{Stealth[length=2.2mm]}, thick, color=secundaria},
      font=\footnotesize
    ]
    \node[box] (user) at (0,0) {Usuário\\(navegador)};
    \node[box] (next) at (4.4,0) {Next.js (App Router)\\Server Components · Actions\\Route Handlers · Auth.js};
    \node[box] (pg) at (8.8,-1.5) {PostgreSQL\\(dados)};
    \node[box] (redis) at (8.8,1.5) {Redis · BullMQ\\(filas)};
    \node[box] (worker) at (12.6,1.5) {Worker\\(e-mails)};
    \draw[seta] (user) -- (next);
    \draw[seta] (next) -- node[above] {Prisma} (pg);
    \draw[seta] (next) -- node[above] {fila} (redis);
    \draw[seta] (redis) -- node[above] {consome} (worker);
    \draw[seta, dashed] (pg.north) -- node[below] {resposta} (next.south east);
  \end{tikzpicture}}
  \caption{Arquitetura do SkillHub: uma aplicação, dados síncronos (banco) e
  trabalho assíncrono (filas).}
  \label{fig:skillhub-arquitetura}
\end{figure}

\begin{table}[h]
\centering
\small
\caption{As três camadas de escrita do App Router e quando usar cada uma.}
\label{tab:camadas-next}
\begin{tabularx}{\textwidth}{@{}>{\raggedright\arraybackslash}p{2.6cm} >{\raggedright\arraybackslash}X >{\raggedright\arraybackslash}X >{\raggedright\arraybackslash}X@{}}
\toprule
\textbf{Pergunta} & \textbf{Server Component} & \textbf{Server Action} & \textbf{Route Handler} \\
\midrule
O que faz? & Lê dados e renderiza & Escreve dados (mutações) & Expõe API HTTP \\
Onde roda? & Servidor (SSR/SSG) & Servidor & Servidor \\
Como é chamado? & Pela renderização & Por um \texttt{<form>} & Por \texttt{fetch} \\
Cliente vê? & Só o HTML & Nada do código & Só a resposta \\
Quando usar? & Listagens, detalhes & Formulários, botões & Webhooks, mobile, terceiros \\
\bottomrule
\end{tabularx}
\end{table}

\begin{caixadica}
Regra prática: \emph{leitura} vai em Server Component, \emph{escrita} vai em
Server Action, e \emph{integração com o mundo externo} vai em Route Handler.
Se você seguir só essa regra, já acerta 90\% dos casos reais.
\end{caixadica}

\section{Buscando dados: do banco à tela}

Um Server Component é uma função \texttt{async} que roda no servidor. Como ele
tem acesso direto ao Node.js, pode importar o Prisma e consultar o banco sem
passar por HTTP:

\begin{lstlisting}[style=livro, language=TypeScript,
  caption={Server Component lendo o banco direto.}, label={list:fullstack-server}]
import { prisma } from "@/lib/prisma";

export default async function PaginaDeServicos() {
  const servicos = await prisma.servico.findMany({
    orderBy: { criadoEm: "desc" },
    take: 20,
  });

  return (
    <ul>
      {servicos.map((s) => (
        <li key={s.id}>
          <h2>{s.titulo}</h2>
          <p>{s.descricao}</p>
        </li>
      ))}
    </ul>
  );
}
\end{lstlisting}

Note o que \emph{não} acontece aqui: não há \texttt{fetch}, não há URL de API,
não há serialização JSON e não há estado de carregamento no cliente. O banco é
consultado durante a renderização no servidor, e o navegador recebe HTML
pronto. Isso elimina o ``cascata de requisições'' típica de SPAs, em que a
página carrega, depois busca dados, depois renderiza.

\begin{caixaconceito}
\textbf{Por que isso é seguro?} O código do Server Component nunca chega ao
navegador — nem o import do Prisma, nem a string de conexão, nem as queries.
Se você tentar importar um módulo de servidor num Client Component, o Next.js
falha em tempo de compilação com um erro claro. Essa fronteira é o que torna o
modelo seguro por construção, e não por disciplina.
\end{caixaconceito}

\subsection{Cache e revalidação}

Por padrão, o Next.js tenta ser o mais estático possível: páginas e dados são
cacheados, e o site fica rápido. O problema: depois de uma mutação, o cache
pode servir dados velhos. A solução oficial é invalidar \emph{sob demanda}:

\begin{itemize}[leftmargin=*, itemsep=0.4em]
  \item \texttt{revalidatePath("/servicos")} — invalida uma rota específica;
  \item \texttt{revalidateTag("servicos")} — invalida todos os dados marcados com aquela tag (\texttt{fetch} com \texttt{next: \{ tags: [...] \}});
  \item \texttt{router.refresh()} — no cliente, recarrega a rota atual sem perder estado.
\end{itemize}

\begin{caixadica}
Use tags quando a mesma informação aparece em várias rotas (por exemplo, a
contagem de pedidos no cabeçalho e no painel). Use \texttt{revalidatePath}
quando a invalidação é local a uma rota. Misturar os dois é comum e saudável.
\end{caixadica}

\section{Mutações com Server Actions}

Uma Server Action é uma função \texttt{async} marcada com \texttt{"use server"}
que o navegador chama diretamente — sem escrever um endpoint, sem
\texttt{fetch}, sem JSON manual. O Next.js gera um POST interno, serializa o
\texttt{FormData} e executa a função no servidor com sua sessão real:

\begin{lstlisting}[style=livro, language=TypeScript,
  caption={Server Action com validação e revalidação.}, label={list:fullstack-action}]
"use server";

import { revalidatePath } from "next/cache";
import { z } from "zod";
import { prisma } from "@/lib/prisma";

const schema = z.object({
  titulo: z.string().min(3).max(80),
  precoCentavos: z.coerce.number().int().positive(),
});

export async function criarServico(prev: Estado, formData: FormData) {
  const parsed = schema.safeParse({
    titulo: formData.get("titulo"),
    precoCentavos: formData.get("precoCentavos"),
  });

  if (!parsed.success) {
    return { erro: "Dados inválidos. Verifique título e preço." };
  }

  await prisma.servico.create({ data: parsed.data });
  revalidatePath("/servicos"); // atualiza a página em cache
  return { sucesso: true };
}
\end{lstlisting}

Três detalhes merecem destaque:

\begin{enumerate}[leftmargin=*, itemsep=0.4em]
  \item \textbf{Validação no servidor}: o schema Zod\index{Zod} roda no servidor. O
  navegador pode enviar qualquer coisa — inclusive dados manipulados — e a
  ação rejeita. Validar só no cliente é enfeite; validar no servidor é segurança.
  \item \textbf{Assinatura com estado}: o primeiro parâmetro \texttt{prev}
  carrega o retorno da execução anterior, permitindo feedback de erro sem
  estado global no cliente.
  \item \textbf{Revalidação}: sem \texttt{revalidatePath}, o usuário publica o
  serviço, mas a listagem em cache continua sem ele — a mudança ``não
  acontece''.
\end{enumerate}

\subsection{O formulário no cliente}

Do lado do cliente, o React 19 entrega \texttt{useActionState}, que liga o
formulário à ação e devolve estado e indicador de progresso:

\begin{lstlisting}[style=livro, language=TypeScript,
  caption={Formulário com useActionState (React 19).}, label={list:fullstack-form}]
"use client";

import { useActionState } from "react";
import { criarServico } from "./acoes";

export function FormularioDeServico() {
  const [estado, acao, pendente] = useActionState(criarServico, {});

  return (
    <form action={acao}>
      <input name="titulo" placeholder="Título do serviço" required />
      <input name="precoCentavos" type="number" placeholder="Preço (centavos)" required />
      <button disabled={pendente}>
        {pendente ? "Salvando..." : "Publicar serviço"}
      </button>
      {estado?.erro && <p role="alert">{estado.erro}</p>}
    </form>
  );
}
\end{lstlisting}

\begin{caixaconceito}
\textbf{Server Actions} transformam formulários HTML em mutações seguras: o
código roda no servidor (nunca exposto ao cliente), o estado de envio vem de
graça via \texttt{useActionState} (\texttt{useFormState} no React 19.1+), e
\texttt{revalidatePath} atualiza o cache. É o padrão mais produtivo do
Next.js 16 — e o que o mercado espera que você saiba.
\end{caixaconceito}

\begin{caixaarmadilha}
\textbf{Não exponha a ação em Client Components arbitrários.} Uma Server
Action importada por um Client Component vira um POST público: qualquer pessoa
pode chamá-la com um \texttt{fetch} — inclusive fora do seu formulário. Por
isso a autorização \emph{precisa} ser revalidada dentro da ação, nunca
assumida pelo fato de o botão estar visível na tela.
\end{caixaarmadilha}

\section{Route Handlers e APIs públicas}

Nem tudo é página. Webhooks de pagamento, aplicativos mobile e integrações com
terceiros precisam de uma API HTTP clássica. Para esses casos, o App Router
oferece \texttt{app/api/.../route.ts} — um arquivo que exporta funções por
método HTTP:

\begin{lstlisting}[style=livro, language=TypeScript,
  caption={Route Handler com validação e status HTTP.}, label={list:fullstack-route}]
import { NextResponse } from "next/server";
import { prisma } from "@/lib/prisma";

export async function GET() {
  const servicos = await prisma.servico.findMany({ take: 50 });
  return NextResponse.json(servicos);
}

export async function POST(req: Request) {
  const body = await req.json();
  // valide com zod, verifique auth, então crie
  return NextResponse.json({ id: novo.id }, { status: 201 });
}
\end{lstlisting}

\begin{itemize}[leftmargin=*, itemsep=0.4em]
  \item \textbf{GET} — leitura pública; o resultado pode ser cacheado com
  \texttt{revalidateTag} no fetch do consumidor;
  \item \textbf{POST/PUT/DELETE} — mutações vindas de fora; valide o corpo,
  autentique o chamador (Bearer token, sessão, assinatura) e retorne códigos
  de status corretos (201, 400, 401, 403, 404);
  \item \textbf{Webhooks} — receba eventos de serviços externos (Stripe,
  GitHub) e processe de forma idempotente: guarde o \texttt{id} do evento para
  não processar o mesmo evento duas vezes.
\end{itemize}

\begin{caixadica}
Regra de ouro: se o consumidor é o \emph{seu próprio frontend}, prefira Server
Components e Server Actions — eles são mais rápidos e mais simples. Se o
consumidor é \emph{outro sistema}, use Route Handlers. A API pública é uma
fronteira; a renderização, não.
\end{caixadica}

\section{Protegendo páginas e ações}

Autenticação e autorização não mudam de natureza no Next.js — elas mudam de
lugar. Tudo acontece no servidor:

\begin{itemize}[leftmargin=*, itemsep=0.4em]
  \item \textbf{Layout protegido}: chame \texttt{auth()} do Auth.js no topo do \texttt{layout.tsx} e redirecione para login;
  \item \textbf{Ações protegidas}: revalide a sessão \emph{dentro} da Server Action (nunca confie no cliente);
  \item \textbf{RBAC}: verifique o papel do usuário logado antes de cada mutação;
  \item \textbf{CSRF}: Server Actions e cookies \texttt{httpOnly} já tratam a maior parte; mantenha-se nos padrões oficiais.
\end{itemize}

\begin{lstlisting}[style=livro, language=TypeScript,
  caption={Autorização dentro de uma Server Action (RBAC).}, label={list:fullstack-rbac}]
"use server";

import { auth } from "@/auth";
import { prisma } from "@/lib/prisma";

export async function excluirServico(id: string) {
  const sessao = await auth();
  if (!sessao?.user) throw new Error("Não autenticado");

  const servico = await prisma.servico.findUnique({ where: { id } });
  if (!servico) throw new Error("Serviço não encontrado");
  if (servico.donoId !== sessao.user.id) {
    throw new Error("Sem permissão para excluir este serviço");
  }

  await prisma.servico.delete({ where: { id } });
  revalidatePath("/servicos");
}
\end{lstlisting}

Repare na sequência: autentica $\rightarrow$ busca o recurso $\rightarrow$
compara o dono $\rightarrow$ age. Se qualquer etapa falhar, a ação falha —
mesmo que o usuário tenha aberto o DevTools e editado o HTML para ``mostrar''
o botão de excluir.

% ============================================================================
\section{Projeto do capítulo: marketplace de serviços}

\begin{caixaprojeto}
\textbf{Objetivo:} construir o \emph{SkillHub} — um marketplace onde
profissionais anunciam serviços (aulas, design, consertos) e clientes fazem
pedidos.

\textbf{Critérios de aceite:}
\begin{itemize}[leftmargin=*, itemsep=0.2em]
  \item Autenticação com Auth.js (e-mail + senha, com hash argon2);
  \item Anúncio de serviços via Server Action (título, descrição, preço, categoria);
  \item Página pública \texttt{/servicos} com busca e filtro por categoria;
  \item Detalhe \texttt{/servicos/[id]} com botão ``Contratar'' que cria um pedido;
  \item Área do profissional \texttt{/painel} listando seus anúncios e pedidos;
  \item RBAC: apenas o dono edita/exclui o anúncio;
  \item Cache invalidado corretamente após cada mutação;
  \item Migrações Prisma + seed com 15 serviços reais.
\end{itemize}
\end{caixaprojeto}

\subsection{Passo a passo}

\begin{passos}
  \item \textbf{Base}: crie o projeto com \texttt{create-next-app} (TypeScript,
  App Router, Tailwind) e configure o Prisma com o PostgreSQL local
  (capítulo~\ref{cap:postgresql});
  \item \textbf{Schema}: modele \texttt{Usuario}, \texttt{Servico} e
  \texttt{Pedido} com relações — um usuário tem vários serviços; um pedido
  aponta para serviço e cliente; rode a primeira migração;
  \item \textbf{Autenticação}: configure o Auth.js com credenciais, hash de
  senha com \texttt{argon2} e sessão via cookie \texttt{httpOnly};
  \item \textbf{Listagem}: implemente \texttt{/servicos} como Server Component
  com busca (\texttt{searchParams}) e filtro por categoria;
  \item \textbf{Criação}: implemente a Server Action \texttt{criarServico}
  (validação Zod + RBAC) e o formulário com \texttt{useActionState};
  \item \textbf{Detalhe e pedido}: \texttt{/servicos/[id]} com Server
  Component, e a ação \texttt{contratar} que cria um \texttt{Pedido} e
  revalida o painel;
  \item \textbf{Painel}: \texttt{/painel} com os anúncios e pedidos do usuário
  logado (layout protegido com \texttt{auth()});
  \item \textbf{Seed}: escreva um script que cria 15 serviços realistas em
  categorias variadas e valide as páginas de ponta a ponta.
\end{passos}

\begin{caixadica}
O SkillHub é o ``capítulo-síntese'' da parte prática: ele reúne 60\% de tudo
o que você aprendeu. Guarde-o — ele vira o case central do seu portfólio.
\end{caixadica}

\section{Exercícios de fixação}

\begin{exercicios}
  \item Quando usar Server Component, Server Action e Route Handler? Dê um caso para cada.
  \item O que \texttt{revalidatePath} faz e por que é necessário após uma mutação?
  \item Escreva uma Server Action que cria um comentário e valida com Zod.
  \item Por que a sessão deve ser revalidada \emph{dentro} da Server Action?
  \item Monte o formulário de criação com \texttt{useActionState} mostrando estado de erro e de carregamento.
\end{exercicios}

\section{Desafios}

\begin{desafios}
  \item \textbf{Upload de imagem}: use \texttt{uploadthing} ou \texttt{s3} para permitir foto no anúncio (com validação de tipo e tamanho).
  \item \textbf{Comentários e avaliações}: adicione comentários com nota (1--5) por serviço, com média exibida.
  \item \textbf{Busca com filtros}: combine categoria + faixa de preço + ordenação na listagem.
\end{desafios}

\section{Exercícios de nível sênior}

\begin{seniores}
  \item \textbf{Otimismo e concorrência}: use \texttt{useOptimistic} para o botão ``Contratar'' atualizar a UI instantaneamente e reconciliar com o servidor.
  \item \textbf{Streaming de páginas}: divida a página de listagem em seções com \texttt{Suspense} e meça a diferença de LCP.
  \item \textbf{Auditoria de cache}: explique a interação entre \texttt{revalidatePath}, \texttt{revalidateTag} e o cache de dados do Next 16 em um diagrama.
  \item \textbf{Testes integrados}: escreva testes E2E do fluxo ``cadastrar $\rightarrow$ anunciar $\rightarrow$ contratar'' (configuração no capítulo~\ref{cap:testes}).
\end{seniores}

\section{Armadilhas comuns}

\begin{itemize}[leftmargin=*, itemsep=0.4em]
  \item Colocar toda a página em \texttt{"use client"} (perde o modelo do servidor);
  \item Validar apenas no cliente (dados podem chegar manipulados);
  \item Esquecer \texttt{revalidatePath} e o usuário não ver a mudança;
  \item Confiar no cliente para autorização (todo request pode ser forjado);
  \item Ações sem tratamento de erro (deixe o usuário saber o que falhou);
  \item Criar \texttt{route.ts} para o próprio frontend (duplica lógica e adiciona latência sem ganho);
  \item Usar \texttt{fetch} dentro de Server Components para chamar a própria API (o Next.js pode otimizar, mas consultar o banco direto é mais simples e rápido).
\end{itemize}

\section{Resumo do capítulo}

\begin{itemize}[leftmargin=*, itemsep=0.3em]
  \item Next.js integra frontend, backend e banco em um único projeto;
  \item Server Components leem o banco; Server Actions escrevem; Route Handlers expõem APIs;
  \item \texttt{revalidatePath/Tag} mantêm o cache coerente;
  \item Valide no servidor e revalide a sessão em toda mutação;
  \item Projeto entregue: marketplace SkillHub completo.
\end{itemize}

\section{Referências oficiais para aprofundar}

\begin{itemize}[leftmargin=*, itemsep=0.3em]
  \item Next.js — Server Actions: \url{https://nextjs.org/docs/app/building-your-application/data-fetching/server-actions-and-mutations}
  \item Next.js — Route Handlers: \url{https://nextjs.org/docs/app/building-your-application/routing/route-handlers}
  \item Next.js — Caching: \url{https://nextjs.org/docs/app/building-your-application/caching}
  \item React — \texttt{useActionState}: \url{https://react.dev/reference/react/useActionState}
\end{itemize}
